\begin{tikzpicture}[
  font=\small,
  node distance=8mm and 12mm,
  box/.style={
    rectangle, rounded corners=2pt,
    draw, very thick,
    align=left,
    inner xsep=6pt, inner ysep=5pt,
    text width=3.5cm
  },
  title/.style={font=\bfseries, align=left, text width=5.3cm},
  midbox/.style={
    rectangle, rounded corners=2pt,
    draw, very thick,
    align=left,
    inner xsep=7pt, inner ysep=6pt,
    text width=3.5cm
  },
  note/.style={font=\footnotesize\itshape, inner sep=1pt},
  arrIn/.style={-Latex, thick, dashed},
  arrOut/.style={-Latex, thick, dotted},
  gapline/.style={line width=2.2pt, draw=gray!65}
]



% ------------------------
% LEFT COLUMN (latent)
% ------------------------
\node[title] (Ltitle) {Decision-relevant (latent) quantities};

\node[box, below=3mm of Ltitle] (L1) {\textbf{Degradation state}\\ \textit{(latent / unobservable)}};
\node[box, below=of L1] (L2) {\textbf{Failure risk / probability}\\ \textit{(decision target)}};
\node[box, below=of L2] (L3) {\textbf{Load--degradation relationship}\\ \textit{(usage $\rightarrow$ wear)}};
\node[box, below=of L3] (L4) {\textbf{Remaining useful life (RUL)}\\ \textit{(conceptual, not directly measured)}};

% ------------------------
% Information gap (structural asymmetry)
% ------------------------
\coordinate (gapX) at ($(L1.east)+(14mm,0)$);

\draw[gapline]
  ($(gapX |- Ltitle.north)+(0,2mm)$) --
  ($(gapX |- L4.south)+(0,-2mm)$);

\node[note, rotate=90, text=gray!70]
  at ($(gapX)+(3mm,10mm )$) {information gap};


% ------------------------
% RIGHT COLUMN (observable)
% ------------------------
\node[title, right=72mm of Ltitle] (Rtitle) {Available and observable\\information};

\node[box, below=3mm of Rtitle] (R1) {\textbf{Proxy usage indicators}\\ \textit{(traffic, schedules, tickets, etc.)}};
\node[box, below=of R1] (R2) {\textbf{Environmental / context variables}\\ \textit{(location, time, season, etc.)}};
\node[box, below=of R2] (R3) {\textbf{Maintenance \& service records}\\ \textit{(delayed, inconsistent, qualitative)}};
\node[box, below=of R3] (R4) {\textbf{Failure events}\\ \textit{(rare; censored by maintenance actions)}};

% ------------------------
% MIDDLE (inference)
% ------------------------
\coordinate (midTop) at ($(Ltitle.east)!0.5!(Rtitle.west)$);
\node[midbox, below=12mm of midTop] (M) {
  \textbf{Probabilistic inference under asymmetry}\vspace{1mm}
  \begin{itemize}\setlength{\itemsep}{0.5mm}\setlength{\topsep}{0mm}
    \item latent state estimation
    \item uncertainty propagation
    \item population-level aggregation
  \end{itemize}
};

% Anchor points on the middle box for cleaner routing (avoid stacked arrows)
\coordinate (Me1) at ($(M.east)+(0,11mm)$);
\coordinate (Me2) at ($(M.east)+(0,4mm)$);
\coordinate (Me3) at ($(M.east)+(0,-3mm)$);
\coordinate (Me4) at ($(M.east)+(0,-20mm)$);

\coordinate (Mw1) at ($(M.west)+(0,10mm)$);
\coordinate (Mw2) at ($(M.west)+(0,3mm)$);
\coordinate (Mw3) at ($(M.west)+(0,-4mm)$);
\coordinate (Mw4) at ($(M.west)+(0,-20mm)$);

% ------------------------
% Incoming links (observable -> inference)
% Put labels ON arrows; curve slightly to avoid riding box borders
% ------------------------
\draw[arrIn] (R1.west) to[out=180,in=0] node[midway, above, note] {indirect} (Me1);
\draw[arrIn] (R2.west) to[out=180,in=0] node[midway, above, note] {contextual} (Me2);
\draw[arrIn] (R3.west) to[out=180,in=0] node[midway, above, note] {delayed} (Me3);
\draw[arrIn] (R4.west) to[out=180,in=0] node[midway, above, note] {rare / censored} (Me4);

% ------------------------
% Outgoing links (inference -> latent decision quantities)
% Labels also on arrows
% ------------------------
\draw[arrOut] (Mw1) to[out=180,in=0] node[midway, above, note] {uncertain} (L1.east);
\draw[arrOut] (Mw2) to[out=180,in=0] node[midway, above, note] {probabilistic} (L2.east);
\draw[arrOut] (Mw3) to[out=180,in=0] node[midway, above, note] {model-based} (L3.east);
\draw[arrOut] (Mw4) to[out=180,in=0] node[midway, above, note] {distributional} (L4.east);



\end{tikzpicture}